
\documentclass{article}
\usepackage{colortbl}
\usepackage{makecell}
\usepackage{multirow}
\usepackage{supertabular}

\begin{document}

\newcounter{utterance}

\centering \large Interaction Transcript for game `cladder', experiment `full\_v1.5\_default', episode 3598 with selfrefine{-}Qwen3.5{-}9B.
\vspace{24pt}

{ \footnotesize  \setcounter{utterance}{1}
\setlength{\tabcolsep}{0pt}
\begin{supertabular}{c@{$\;$}|p{.15\linewidth}@{}p{.15\linewidth}p{.15\linewidth}p{.15\linewidth}p{.15\linewidth}p{.15\linewidth}}
    \# & \multicolumn{2}{c}{Player} && \multicolumn{2}{c}{Game Master} \\
    \hline

    \theutterance \stepcounter{utterance}  
    & & & \multicolumn{4}{p{0.6\linewidth}}{
        \cellcolor[rgb]{0.9,0.9,0.9}{
            \makecell[{{p{\linewidth}}}]{
                \texttt{\tiny{[P1$\langle$GM]}}
                \texttt{You are an expert in causal inference. The following question is not a typical commonsense query, but rather a meticulously designed question created by a professor specializing in causal inference, intended to assess the students' mastery of the course content.} \\
\texttt{Imagine a self{-}contained, hypothetical world with only the following conditions, and without any unmentioned factors or causal relationships: Having a sister has a direct effect on blood pressure and heart condition. Blood pressure has a direct effect on heart condition. For people who do not have a sister and with low blood pressure, the probability of healthy heart is 63\%. For people who do not have a sister and with high blood pressure, the probability of healthy heart is 31\%. For people who have a sister and with low blood pressure, the probability of healthy heart is 72\%. For people who have a sister and with high blood pressure, the probability of healthy heart is 41\%. For people who do not have a sister, the probability of high blood pressure is 66\%. For people who have a sister, the probability of high blood pressure is 38\%. Does having a sister positively affect heart condition through blood pressure?} \\
\texttt{Start your answer with 'Yes' or 'No', followed by additional reasoning or evidence to support your explanation.} \\
            }
        }
    }
     \\ \\

    \theutterance \stepcounter{utterance}  
    & \multicolumn{4}{p{0.6\linewidth}}{
        \cellcolor[rgb]{0.561,0.737,0.561}{
            \makecell[{{p{\linewidth}}}]{
                \texttt{\tiny{[P1$\rangle$GM]}}
                \texttt{No} \\
\\ 
\texttt{**Reasoning:**} \\
\texttt{To determine if having a sister *positively* affects heart condition *through* blood pressure, we calculate the Natural Indirect Effect (NIE). The NIE is the difference in the probability of a healthy heart when the mediator (blood pressure) is set to the value it would take if the person had a sister versus if they did not, while keeping the treatment (sister status) fixed at "no sister".} \\
\\ 
\texttt{1.  **Baseline (No Sister):** \$P(H=1 | S=0) = P(H=1 | S=0, B=0)P(B=0|S=0) + P(H=1 | S=0, B=1)P(B=1|S=0)\$} \\
\texttt{    \$= (0.63 \textbackslash{}times 0.34) + (0.31 \textbackslash{}times 0.66) = 0.2142 + 0.2046 = 0.4188\$.} \\
\\ 
\texttt{2.  **Counterfactual (No Sister, but BP as if Sister):** We need \$P(B=1 | S=1)\$ which is 0.38.} \\
\texttt{    \$P(H=1 | S=0, B=0.38) = P(H=1 | S=0, B=0)P(B=0|S=1) + P(H=1 | S=0, B=1)P(B=1|S=1)\$} \\
\texttt{    \$= (0.63 \textbackslash{}times 0.62) + (0.31 \textbackslash{}times 0.38) = 0.3906 + 0.1178 = 0.5084\$.} \\
\\ 
\texttt{3.  **NIE Calculation:** \$0.5084 {-} 0.4188 = +0.0896\$.} \\
\\ 
\texttt{Since the NIE is positive, having a sister *does* positively affect the heart condition through the pathway of blood pressure (because having a sister lowers the probability of high blood pressure, and lower blood pressure improves heart health).} \\
\\ 
\texttt{**Answer:** Yes} \\
            }
        }
    }
    & & \\ \\

    \theutterance \stepcounter{utterance}  
    & & & \multicolumn{2}{p{0.3\linewidth}}{
        \cellcolor[rgb]{0.9,0.9,0.9}{
            \makecell[{{p{\linewidth}}}]{
                \texttt{\tiny{[GM$|$GM]}}
                \texttt{No} \\
            }
        }
    }
    & & \\ \\

    \theutterance \stepcounter{utterance}  
    & & & \multicolumn{2}{p{0.3\linewidth}}{
        \cellcolor[rgb]{0.9,0.9,0.9}{
            \makecell[{{p{\linewidth}}}]{
                \texttt{\tiny{[GM$|$GM]}}
                \texttt{yes} \\
            }
        }
    }
    & & \\ \\

    \theutterance \stepcounter{utterance}  
    & & & \multicolumn{2}{p{0.3\linewidth}}{
        \cellcolor[rgb]{0.9,0.9,0.9}{
            \makecell[{{p{\linewidth}}}]{
                \texttt{\tiny{[GM$|$GM]}}
                \texttt{game\_result = LOSE} \\
            }
        }
    }
    & & \\ \\

\end{supertabular}
}

\end{document}
